% =====================================================================
% CRITICAL EVALUATION (part f)
% LLM declaration: text level 3; see the appendix "Use of AI/LLM tools".
% =====================================================================

\section*{Critical evaluation}

\emph{Regression.} OLS has a closed-form solution. It takes milliseconds, gives the same result every time and needs no learning rate, and with a good degree it was as good as our first networks. But we must choose the features ourselves, a polynomial of high degree oscillates near the edges and blows up outside the data (Fig.~\ref{fig:b_fits}), and the best degree depends on the data set. Ridge keeps the high degrees under control with a small penalty, still in closed form, and was the best method in Project 1. Lasso has no closed form, and its solver was slow and did not always converge. The network learns its features from $x$ alone, and the ReLU network gave the lowest test error of all the methods, close to the noise level, mainly by doing better near the edges. On the other hand it has many hyperparameters (architecture, activation, optimiser, learning rate, minibatch size, number of epochs and penalty), its cost is not convex, so the result depends on the initial weights and the minibatches, and it trains in seconds instead of milliseconds. With only 16 validation points the choice of hyperparameters was noisy, and $k$-fold cross-validation on the 80 training points, as in Project 1, would have used the data better. For a smooth function of one variable, OLS or Ridge is the natural first choice. The network pays off when the function is hard to describe with a polynomial, as near the edges here.

\emph{Classification.} Logistic regression has a convex cost with one minimum, few hyperparameters and short training, and its weights can be read as templates. It is the baseline a network must beat. The network with one hidden layer beat it clearly, with 69\,\% fewer errors, because the features that make up a handwritten digit are not known in advance, and the hidden layer, with its non-linear activation, learns them. This is where the network earns its extra cost, about ten times the training time and more hyperparameters. Of the methods we tried, it is the best for the classification; convolutional networks do better still.

\emph{Gradient methods.} Plain gradient descent has only one parameter, $\gamma$, but its step is proportional to the gradient. It needed a learning rate 100 times larger than RMSprop and Adam, overshot at the start and converged slowly, and its best $\gamma$ changes with the architecture. RMSprop and Adam scale each step to about $\gamma$, so $\gamma = 10^{-3}$ worked well for all the architectures and on both data sets, and with minibatches they make several updates per epoch. Adam was the least sensitive to $\gamma$, and we used it for parts c--e. With a constant learning rate and a fixed number of epochs, when to stop matters: trained longer, the networks started to follow the noise, and on MNIST the accuracy jumped from epoch to epoch.
